\section{Why each ingredient is needed}\label{sec:ingredients}

Beyond the label framework of the eighth paper~\cite{EighthPaper2026}, the proof uses three ingredients: bands measured from
homes adapted to the walk (the ladders and frames of \cref{sec:ladders}, with their hybrid running maximum), displays of
suspended frames (\cref{sec:displays}), and random break bits. This section records why simpler choices fail. The remarks are
not formalized and no proof in this paper uses them; each is labelled as a proof sketch, a heuristic or numerical.

\begin{remark}[Homes at turning points fail]\label{rem:runstart}
\emph{Proof sketch only.} The simplest band measures the distance from the start of the current monotone run of the walk: it is
the band of the frame system whose only frame is the root and whose restart points are the times at which the walk reverses
direction. Theorem~A fails for it on the \emph{nested mountains}. Let $\mathrm{St}(g)=(+1,+1,-1)^g$ be a slow staircase and
$\mathrm{Dn}(g)=(-1)^g$ a fast descent, put $M_0=(+1,+1,-1,-1)$ and, for $j\ge1$ and integers $0\le a_j\le2^{j-1}$,
\[
  M_j=\mathrm{St}(a_j)\,M_{j-1}\,M_{j-1}\,\mathrm{St}(2^j-a_j)\,\mathrm{Dn}(2^j),
\]
and take the walk with the steps of $M_J$, of length $T=4(J+1)2^J$. A copy of~$M_j$ climbs slowly by~$2^j$, carrying two copies
of $M_{j-1}$ that stay below its summit, and returns in a single run. On a slow climb the last visit to a cell is a reversal, so
the hops that end on the lower half of the final descent of a copy of~$M_j$ start at band~$0$; they end at distance more than
$2^{j-1}$ from the summit, at band at least $j-1$, and their extent exceeds $2^{j-1}$. For $j\ge h_0+1$ their two ends therefore
lie in different groups at every offset, and by (P1) of \cref{prop:prom} their owners have scale at least $j-1\ge h_0$, so they
are not ground: they are lost (\cref{def:kept}) at every offset. Each~$j$ contributes at least $2^{J-j}\cdot2^{j-1}=T/(8(J+1))$
such hops, so at least $(J-h_0)\,T/(8(J+1))$ hops are lost at every offset, a fraction of~$T$ that tends to~$\frac18$ as $J$
grows, whereas Theorem~A allows $O(T/h_0)$. Bands that measure from turning points measure speed, not distance. Measuring the
distance from the base of the excursion instead (\cref{sec:ladders}) makes these walks cheap: the climb of an excursion has no
restart, and its descent restarts only at the passages $D(y+2^e)$.
\end{remark}

\begin{remark}[Frames must be displayed]\label{rem:popfreeze}
\emph{Heuristic.} Frames make the hops of each walk cheap on their own (\cref{thm:A}), but they create an interaction between
walks. While a frame runs, the segment that it suspends keeps its band frozen; when the segment resumes, its band group
reappears. A label that shows only the current band groups of the walks therefore changes its shown set at every pop at a high
band. Such a change can produce a descent (\cref{lem:stateedge}), and it falls under the case (ABOVE) of Lemma~V (\cref{lem:V})
for every hop of another walk that spans it with a lower deciding group. On walks whose long excursions carry dense nested
structures at many scales these pops are frequent, and the label pays for them. Either the frozen band is shown while the frame
runs, or the pop that ends the frame changes the shown set at a high group; this pop-or-freeze dichotomy forces the label to show
suspended frames, which is what the entries of \cref{sec:displays} do.
\end{remark}

\begin{remark}[Exact displays are unbounded]\label{rem:exact}
\emph{Proof sketch only.} Call a display \emph{exact} if, with every entry unbroken, its shown set above every group~$d$ depends
only on the view above~$d$ (\cref{sec:displays-view}), so that view invariance alone, with no random bits, would control every
hop. The view above~$d$ does not distinguish a candidate at~$d$ from a seed at its top, and seeds are always shown; so an exact
display shows the top of every candidate at every value, and in particular the whole cascade of windows: an entry whose window
contains a shown value is shown, and so on upwards. On the \emph{dense tower}, a walk whose stack of suspended frames carries,
at a constant fraction of times, unboundedly many entries with distinct tops and windows of more than one group, exact displays
therefore have unbounded size, and no bounded number of digits can carry them. Random break bits cut the cascade at the first
broken entry: Lemma~D (\cref{lem:D}) bounds the size of the display independently of the depth of the stacks, and a hop is
exposed only through a broken candidate, with probability at most $2n/Q$ (\cref{prop:exposure}).
\end{remark}

\begin{remark}[The three definitional choices]\label{rem:fixes}
\emph{Numerical.} Three details of \cref{def:entries,def:view} matter: an entry is a suspend child (its parent segment is
innermost again after it ends), its top is the group of the frozen band of its parent segment at the push, and a rise of the
running maximum at the time a frame resumes is a standard event. Under each alternative reading (every suspend frame is an
entry; the top is the head's band group just before the push, which differs from the frozen band of the parent at a
simultaneous push; a rise at a resumption is not an event), an exact search over pairs of short walks, with $h_0=1$ and all
bits unbroken, finds a kept hop whose shown sets above its deciding group differ at its two ends although no standard event
lies between them; the smallest instances found have lengths $23$, $41$ and~$23$. So the analogue of view invariance fails
under each alternative reading. With the three choices, view invariance holds for all walks (\cref{prop:W}).
\end{remark}
